\summarypart{TWO}{The Age of Catholic Christianity}{70--312}
\begin{partoverview}
Between the fall of Jerusalem and the rise of Constantine the churches multiplied in every province of
the Roman Empire and beyond its eastern border, perhaps as far as India. Their members thought of
themselves as one body and used the word \emph{catholic} for it. The term carried a double claim.
Against mockery and persecution it asserted that the faith was meant for the whole world; against rival
versions of the gospel it asserted that this was the genuine teaching of Jesus. Under these pressures
the bishop became the central figure of each church. The part therefore follows three developments: a
mission without frontiers, an agreed body of belief, and leadership by bishops.
\end{partoverview}
\partchapters{%
  \pch{3}{Only Worthless People}{Catholic Christianity}
  \pch{4}{If the Tiber Floods}{The Persecution of Christians}
  \pch{5}{Arguing About the Event}{The Rise of Orthodoxy}
  \pch{6}{The Rule of Books}{The Formation of the Bible}
  \pch{7}{The School for Sinners}{The Power of Bishops}
  \pch{8}{Apostles to Intellectuals}{The Alexandrians}}
